\subsection{Quantitative simulation on biased
products}\label{exch:sec-2}

\textbf{Lemma.} For \(a\in(0,1/2]\), every bias \(q\in[a,1-a]\), every
input length \(N\), and every \(T\)-query acceptance function, there is
a deterministic simulator with error \(e\) under
\(\operatorname{Bernoulli}(q)^{\otimes N}\) and at most

\[
S(T,e,a)=S_0(2bT,2/e),\qquad
b=\max\left\{1,\left\lceil\log_2\frac{64T^2}{ae^2}\right\rceil\right\}
\]

queries. A monotone polynomial majorant gives a common bound for all
indicated biases.

\textbf{Proof.} Let \(q'\) be the smallest dyadic of denominator \(2^b\)
at least \(q\). A threshold gadget \(g\) on \(b\) uniform bits has
acceptance probability \(q'\). Lift the original quantum algorithm to
\(Q(g(z_1),\ldots,g(z_N))\). Each virtual query is implementable with
\(2b\) input queries, so uniform AA supplies a classical approximation
with error \(e/2\) and the stated query bound.

On a \(q'\)-biased input \(Y\), independently generate each block
conditional on \(g(z_i)=Y_i\) and simulate the classical tree. Formally,
pre-sample a separate conditional block for each possible bit value at
every position, independently of the input, then select the appropriate
block after querying that bit. Lazy generation needs at most one query
to \(Y_i\) when that block is first touched. The expected error is still
\(e/2\); fix the auxiliary randomness by averaging to obtain a
deterministic tree \(A(Y)\).

Couple \(q\)-biased \(X\) and \(q'\)-biased \(Y\) coordinatewise by
common independent uniform thresholds. Then

\[
\Pr[X_i\ne Y_i\mid X_i]\le(q'-q)/(1-q)\le2^{-b}/a.
\]

The quantum hybrid estimate gives \(\mathbb E|Q(X)-Q(Y)|\le e/2\), hence
\(\mathbb E|A(Y)-Q(X)|\le e\). To implement this approximation on \(X\),
sample the independent channels \(Y_i\mid X_i\) as each location is
queried. Fix the channel randomness by averaging. The resulting
deterministic tree on \(X\) uses no additional queries and has error at
most \(e\). All fixed random seeds are input-independent. \(\square\)
